Deploying constrained learning algorithms requires a highly efficient method for computing Euclidean projections. Dykstra's algorithm~\cite{DYKSTRA} offers a computationally lightweight solution to find the projection of an initial point onto the intersection of multiple convex sets~\cite{KEMPF20206542}. However, despite its strong theoretical guarantees, Dykstra's algorithm is mathematically encumbered by the stalling phenomenon~\cite{DYKSTRASTALLING}. Because the duration of these stalling periods is unknown \emph{a priori}, the algorithm cannot be reliably restricted to a fixed iteration budget. This unpredictable convergence behaviour inhibits the deployment of Dykstra's method in practical, real-time computational software.

In this chapter, we address the stalling problem of Dykstra's method for polyhedral sets. Exploiting simplifications that arise in the polyhedral case, we derive the exact closed-form duration of the stalling period once stalling is detected. This mathematical derivation enables a modified version of Dykstra's algorithm that instantly fast-forwards through the stalling period, thereby eliminating unpredictable computational cycles. Crucially, this modification preserves all asymptotic convergence guarantees of the original algorithm. This theoretical development forms the fundamental optimisation pillar of this thesis, transforming Dykstra's method into a robust, predictable engine ready for safety-critical engineering applications.

% The structure of this chapter is as follows. Section~\ref{sec:euclidean_projections} reviews the mathematical background of Euclidean projections and the method of alternating projections. Section~\ref{sec:dykstra_algorithm} introduces Dykstra's algorithm and its specific simplifications for polyhedral sets. Section~\ref{sec:stalling_phenomenon} formally defines the stalling phenomenon. Section~\ref{sec:closed_form_solution} presents the rigorous mathematical proof for the closed-form duration of the stall. Section~\ref{sec:fast_forward_algorithm} proposes the fast-forward algorithmic modification, and Section~\ref{sec:dykstra_numerical_experiments} presents supporting numerical experiments to demonstrate the algorithmic superiority of the proposed framework on pathological constraint geometries.

\section*{Notation}
We denote the transpose of a vector $w$ as $w^T$, and its $\ell_{2}$-norm as $\twonorm{w}$. The interior of a set $\mathcal{H}$ is denoted as $\text{int}\,\mathcal{H}$, and its boundary is denoted as $\partial\mathcal{H}$. Function composition is written as $f \circ g$, and the modulo and ceiling operators are denoted by $[\cdot]$ and $\lceil \cdot \rceil$, respectively. The absolute inner Dykstra iteration index is denoted as $m$ and raised as a bracketed superscript (e.g., $w^{(m)}$), while the outer learning iteration index is denoted as $t$. The spatial dimension of the decision variable is denoted as $d$. The individual half-space index is denoted as $i$, while the total half-space count is denoted as $n$ and remains fixed. We name an update with an increase by $1$ in $m$ ($m \leftarrow m+1$) an \emph{iteration}, and refer to an advancement of $n$ iterations ($m \leftarrow m+n$) as a \emph{cycle}.